% Valorisation strategy - IOF StarTT (AEGIS).
% A4, 10.5pt, sans (Panno substituted by helvet). Max 2 pages.
\documentclass[10pt,a4paper]{article}
\usepackage[T1]{fontenc}
\usepackage[utf8]{inputenc}
\usepackage[margin=2.0cm,top=1.5cm,bottom=1.5cm]{geometry}
\usepackage{helvet}
\renewcommand{\familydefault}{\sfdefault}
\usepackage{booktabs}
\usepackage{microtype}
\usepackage{enumitem}
\usepackage[dvipsnames]{xcolor}
\usepackage{titlesec}
\usepackage{parskip}

\makeatletter
\renewcommand\normalsize{\@setfontsize\normalsize{10.5}{11.6}}
\makeatother
\normalsize

\definecolor{ugblue}{HTML}{1E64C8}
\titlespacing*{\section}{0pt}{2pt}{2pt}
\titleformat{\section}{\bfseries\color{ugblue}\normalsize}{}{0pt}{}
\setlength{\parskip}{1.5pt}
\pagestyle{empty}
\setlist[itemize]{leftmargin=1.1em,itemsep=1pt,topsep=1pt,parsep=0pt}

\begin{document}

\section{Valorisation strategy}

\textbf{First market: absorbed power density for devices above 6\,GHz.} ICNIRP 2020 made absorbed power density the basic
restriction above 6\,GHz, and the IEC/IEEE 63195 series makes it the quantity a millimetre-wave device must be certified
against. Buyers are device makers and their chipset suppliers, the test laboratories that certify them, and the dosimetry
institutes serving both. Two facts define the opening.

\textbf{1. The procedure does not exist yet, and it is being written now.} Parts 1 and 2 of IEC/IEEE 63195 (\emph{incident} power
density) were published in 2022. Parts 3 and 4 (\emph{absorbed}) are still drafts, and part 4, ``Computational Procedures for Absorbed Power Density'', currently specifies FDTD and FEM. The computational
rules for the quantity this project computes are open, in a working group, inside this project's window.
ICNIRP asks for exactly this work in writing: ``information from a technical standards body [\ldots] to more adequately
match the basic restrictions, should be utilized to improve reference level assessment procedures.'' The generic reference levels are a deliberately conservative plane-wave
proxy, and no source-specific procedure replaces them above 6\,GHz.

\textbf{2. The obstacle is cost, and it is ours to remove.} Computing absorbed power density on an anatomical body needs a
full-wave simulation, at hours per scenario. A device product cycle needs thousands: beam configuration by grip by body
position by antenna placement. Nobody can afford that, so the industry falls back on the generic proxy and designs with
margin it cannot quantify. \textbf{AEGIS replaces the volume simulation with a closed-form surface integral: milliseconds per
scenario, and differentiable, so absorbed power becomes a design variable rather than a post-hoc check.} Emissions teams already pay \$500--2{,}000 for a pre-compliance
day to avoid a failed campaign costing \$5{,}000--30{,}000 and four to twelve weeks. RF exposure has no rentable bench, so
pre-compliance here is simulation or nothing.

\textbf{On the size of the market, plainly.} This is a specialist market and I will not dress it up. The
electromagnetic-simulation market is roughly \$1.7B, but the slice concerned
with human RF exposure is a low-single-digit-millions niche, and the figure circulated for ZMT (Sim4Life) is one plausible
order-of-magnitude marker for it. Bottom-up, the buyers are around six infrastructure vendors, twelve to fifteen millimetre-wave
device makers plus their chipset suppliers, twenty to forty RF-exposure test laboratories, ten to fifteen regulators, and several
dozen research institutes. Wi-Fi 6E/7 widens the device tier: every portable computer above 5.925\,GHz now certifies
against power density, and those filings already report absorbed power density routinely (one 2023 filing:
9.85 against the 10\,W/m$^2$ limit), multiplying the device-maker universe at a lower price point, reached
through the laboratories. A serviceable market of a few million
euros a year is the honest ceiling. That is the right size for
this instrument: not a venture-scale company, but a capital-efficient software spin-off on the model of two recent StarTT spin-offs
in this group, targeting seven-figure recurring revenue and a strategic exit to an established simulation vendor, compliance-tool
vendor, or testing-and-certification group, all established acquirers of compliance software.
A small niche is not an argument against the spin-off. It is an argument for reaching it through the incumbents' distribution rather
than around it.

\textbf{Customer discovery: what the buyers actually said.} No company is a project partner. These are prospective customers.
\begin{itemize}
  \item \textbf{CNR-IEIIT Milan, Dr.~M.~Parazzini, Director of Research.} Two of her researchers work above 6\,GHz, more as FR3
    arrives. Her group loses 30 to 50\% of its time waiting for FDTD, and \textit{``the bottleneck is
    always the mesh, it is our nightmare''}. On price: \textit{``if I really find a real improvement with respect to what I have
    now, I would be able to pay you.''} She has drafted a letter of support and licence interest and ranked her adoption needs: usable
    platform with tutorials (5/5), validation paper (4/5), standards recognition (3/5); Phase~1 delivers all three. She is the
    validation and citation route, not the revenue route, and is priced accordingly.
  \item \textbf{Ericsson Research, EMF team} (6--7 people, 1--2 FTE on full-body simulation in CST, FEKO and HFSS). They set the
    target: \textit{``if you can reduce the compliance boundary only because you added the human body, that would be a great
    story.''} They handed over the deal shape unprompted: absorbed power density needs the same input data as the
    incident power density they already compute, so the method is \textit{``a drop-in replacement, for EMF Visual or IXUS.''} Their 63195 contributor added: go to devices, not base stations. Outcome:
    an invitation to demonstrate to the full team after the summer.
  \item \textbf{Nokia, C.~Grangeat}, convenor of IEC TC106 MT3. \textbf{A clear negative, and the most useful result obtained.} He
    reports no operator pull for body-based compliance (\textit{``operators probably prefer on reference level basis''}) and judged
    the base-station route closer to research than to valorisation. I take that at face value: hence the device-side
    beachhead, where absorbed power density is already what must be certified, with base stations scoped to one escalation case.
\end{itemize}

\textbf{Pricing, and the two datapoints that constrain it.} Nokia's EMF team spends under 10\,kEUR a year on software of this
class, and CNR-IEIIT pays 3\,kEUR per exposure licence. Neither supports a seat sold into an EMF team's existing software line,
and ZMT gives academic groups Sim4Life free, so a paid academic tier would not clear.
The model therefore sells design-loop value to the engineering budget that already pays for full-wave tools, and
distribution to the vendors whose compliance products the customers already run.

\textbf{Academic and research: free}, with a citation requirement, a deliberate funnel. \textbf{Test laboratory and institute seat:
15\,kEUR/yr}, below one commercial full-wave seat. \textbf{Device maker and vendor seat} with API, report automation and the design
loop: \textbf{40\,kEUR/yr}, under half an FTE against a pre-compliance workflow of thousands of scenarios per product cycle.
\textbf{OEM licence into a compliance tool}, the route Ericsson described: \textbf{80\,kEUR/yr plus royalty}. \textbf{Contract research:
50\,kEUR} per project. \textbf{Fixed-price pre-compliance studies: 3--8\,kEUR}, the market rate for a simulation
sweep, first revenue and discovery channel.

\begin{center}\small
\begin{tabular}{@{}lccccr@{}}
\toprule
Year & Lab seats & Device seats & OEM licences & Projects & Turnover (EUR) \\
\midrule
1 & 2 & 1 & 0 & 1 & 120\,000 \\
2 & 3 & 2 & 1 & 1 & 255\,000 \\
3 & 5 & 4 & 1 & 2 & 415\,000 \\
4 & 6 & 5 & 2 & 2 & 550\,000 \\
5 & 8 & 6 & 3 & 3 & 750\,000 \\
\bottomrule
\end{tabular}
\end{center}

\noindent Fourteen direct seats by year 5, against a direct universe of sixty to eighty organisations, with a third of revenue
carried by three OEM licences rather than by a penetration rate no one-person spin-off could deliver. Against a project
investment of order 250\,kEUR, that is roughly 2.1\,MEUR cumulative on software margins with no hardware to finance.

\textbf{Route to market: the standard is the channel.} Ericsson said it unprompted: \textit{``it has
to be in the standard, that is always most important for the business side, and for regulators.''} The incumbent proves it: ZMT's position
rests less on its solver than on a computational tool written into a regulatory pathway. That is an
unusually favourable channel here: the promotor is the Belgian representative to CENELEC TC106X and IEC TC106, and WAVES sits in
three of the four EU 5G-health projects. The method is contributed as fast pre-screening complementing FDTD, which is how it is genuinely
used. First sales need not wait for it: EU RED is self-declaration, a simulation report in the technical
file already suffices, and ISED Canada publishes codified simulation procedures (RSS-102.IPD.SIM).

\begin{center}\small
\begin{tabular}{@{}llp{5.9cm}p{4.0cm}@{}}
\toprule
Phase & Timing & Activity & Gate \\
\midrule
1. Validate & M1--M6 & Validation against a reference solver \textbf{and against measurement}: a 63195 device benchmark plus
public FCC filings. Validation paper. Public tutorial platform. \textbf{Also settles the base-station question} (below). &
Agreement within the error bound. Go/no-go on the exploration WP and on base stations. \\
\addlinespace
2. Standardise & M6--M18 & Contribution to the open IEC/IEEE 63195-4 draft and to IEC TR 62669 Ed.\,3, via the promotor & Method
recorded as an accepted computational approach \\
\addlinespace
3. Sell & M9--M24 & Full-day demonstration at Ericsson (invited). Paid pilots with device makers and test labs, prospected from FCC
filing analysis. First OEM licence
conversation with a compliance-tool vendor & Two paid pilots, one letter of intent \\
\addlinespace
4. Incorporate & M18--M24 & Spin-off preparation with TechTransfer and istart & Spin-off founded, first licences signed \\
\bottomrule
\end{tabular}
\end{center}

\textbf{The base-station case is a measurement, not a claim.} Above 6\,GHz the restriction that binds at a compliance boundary is
whole-body SAR at 0.08\,W/kg for the general public, and its reference-level proxy of 10\,W/m$^2$ was carried over unchanged from ICNIRP 1998,
set by whole-body resonance near 70\,MHz and never re-derived for the superficial-absorption regime. My estimate is that an anatomical assessment recovers of order 1.5$\times$ in permitted power, hence around 1.2$\times$ in distance, but no
published value exists. Phase~1 computes it on a named IEC 62232 configuration at 3.5 and
26\,GHz and reports the ratio. If it is 1.0 I will say so, and the case closes at month 6 rather than month 18.

\textbf{Adjacent markets: a bounded exploration on a legal deadline.} The engine is not specific to human tissue. It computes, in
closed form and differentiably, how a coherent field partitions into absorbed, scattered and transmitted power on any electrically large
body. Screening applications with UGent TechTransfer produced a falsifiable selection rule: the method adds value where the task is
to \textbf{design a controllable coherent excitation against a quadratic power limit}, and none where it is to \textbf{certify a passive scattering
signature}. Four pass, among them constrained synthesis of reconfigurable intelligent surfaces and automotive in-cabin sensing. A dedicated
work package runs this to a hard gate: \textbf{a signed letter of intent or a paid pilot from a named company
by M6, or the patent family stays narrow}. Priority is filed in 2026, so the M6 gate lands just before the family's scope can last be widened (PCT
deadline, August 2027). Prof.~T.~Dhaene (IDLab, ex-Keysight) advises it. It is
cheap, gated early, and nothing else depends on it.

\textbf{Valorisation route and team.} The primary route is a \textbf{spin-off}, incorporated at the project's end, as in two recent
StarTT projects in this group. R.~Wydaeghe founds it, Prof.~W.~Joseph (promotor) carries the standards route, UGent TechTransfer the IP and incorporation, plus an istart-recruited co-founder. A licence to an
established vendor is an exit, not an entry: the asset is product plus standards position.

\end{document}
